\documentclass[10pt,a4paper]{amsart}
\usepackage[margin=19mm]{geometry}
\usepackage{amsmath,amssymb,mathtools}
\usepackage[scr=rsfs,cal=euler]{mathalfa}
\usepackage{xfrac}
\usepackage[unicode,hidelinks,bookmarks=false]{hyperref}
\newcommand{\ie}{i.e.\ }
\newcommand{\cf}{cf.\ }
\newcommand{\resp}{respectively\ }
\newcommand{\Subsection}{\subsection}
\providecommand{\nobreakdash}{\nobreak\hbox{-}}
\setlength{\emergencystretch}{2em}
\multlinegap0pt
\usepackage{bm}
\usepackage{bbm}
\usepackage{dsfont}
\usepackage{xspace}
\usepackage{cancel}
\usepackage{mathtools}
\usepackage{amsthm}
\makeatletter
\@ifundefined{thm}{\newtheorem{thm}{Thm}}{}
\@ifundefined{lem}{\newtheorem{lem}[thm]{Lemma}}{}
\makeatother
\DeclareMathOperator{\Dpst}{\mathit{D}_{\mathrm{pst}}}
\DeclareMathOperator{\pst}{pst}
\DeclareMathOperator{\ur}{ur}
\newcommand\End{\mathop{\mathrm{End}}\nolimits}
\newcommand\GL{\mathop{\mathrm{GL}}\nolimits}
\newcommand\GSpin{\mathop{\mathrm{GSpin}}\nolimits}
\newcommand\KS{\mathop{\mathrm{KS}}\nolimits}
\newcommand\so{\mathop{\mathfrak{so}}\nolimits}
\newcommand{\Q}{{\mathds Q}}
\title[Arithmetic monodromy of hyper-K\"ahler varieties over $p$-ad]{Arithmetic monodromy of hyper-K\"ahler varieties over $p$-adic fields}
\author{Kazuhiro Ito, Tetsushi Ito, Teruhisa Koshikawa, Teppei Takamatsu, Haitao Zou}
\date{Source: arXiv:2507.13713v1 (CC BY 4.0)}
\begin{document}
\maketitle

\noindent\emph{Source and licence.} Arithmetic monodromy of hyper-K\"ahler varieties over $p$-adic fields. Version arXiv:2507.13713v1 (CC BY 4.0), distributed under the Creative Commons Attribution 4.0 International licence, as recorded in the arXiv licence metadata for this exact version and shown on the abstract page https://arxiv.org/abs/2507.13713v1. Full rights notice: licenses/arxiv-2507.13713-CC-BY-4.0.txt.\par\smallskip

\noindent\textit{Editorial scope.} Counted unit: the Lem whose source label is \texttt{Lemma:nilpotent operator on Kuga-Satake} in the article (source0.tex lines 1797--1806, source label \texttt{Lemma:nilpotent operator on Kuga-Satake}), delivered together with the article's own complete original proof of it (source0.tex lines 1807--1830, one \texttt{proof} environment of 24 lines). No mathematics is written, repaired or re-derived: statement and proof are the article's own text line for line. Required context retained uncounted: Equation:Kuga-Satake embedding full cohomology (lines 1774--1776). Every definition, display and label of those lines is retained unchanged. Omitted from this package and not redistributed: 1--1773, 1777--1796, 1831--2890. Nothing inside a retained excerpt is omitted or altered: every retained body is delivered exactly as the article's lines are written, so no omission receipt of any kind accompanies this package. Citations: the retained excerpts carry no citation command at all, so no bibliography entry of the article is transcribed and no reference is redistributed. Reference adaptation: none was needed and none was made; every reference inside the delivered unit resolves inside it, so no unresolved reference or citation is produced. Printed slips kept literal: no printed word, symbol, hypothesis or number of the counted statement or of its proof was altered. Layout and class: the article's own class is \texttt{amsart} with packages including geometry, inputenc, xcolor, amsmath, amscd, amssymb, amsthm, extarrows, biblatex, hyperref, url, xurl, stmaryrd, dsfont; the wrapper is the lane's pinned \texttt{amsart} editorial wrapper at 10pt on A4, which restates the theorem-like environments the retained text needs and adds the article's own macro definitions that the retained text needs (\texttt{\textbackslash Dpst}, \texttt{\textbackslash End}, \texttt{\textbackslash GL}, \texttt{\textbackslash GSpin}, \texttt{\textbackslash KS}, \texttt{\textbackslash Q}, \texttt{\textbackslash pst}, \texttt{\textbackslash so}, \texttt{\textbackslash ur}), with extra wrapper packages (bm, bbm, dsfont, xspace, cancel, mathtools). The delivered numbering is the unit's own. Assets: no retained span contains a figure environment, a graphics-inclusion command, a PDF-page inclusion, a table or a TikZ picture; no image file is copied into this package, the package declares no asset dependency and no image file is redistributed with it. Licence: version arXiv:2507.13713v1 of this article is distributed under the Creative Commons Attribution 4.0 International licence, as recorded in the arXiv licence metadata for this exact version and shown on the abstract page https://arxiv.org/abs/2507.13713v1. A full rights notice is delivered at licenses/arxiv-2507.13713-CC-BY-4.0.txt. Characters: every retained excerpt is pure ASCII, so the delivered engine, XeLaTeX with its default Latin Modern text font, reports no missing character.
\begin{equation}\label{Equation:Kuga-Satake embedding full cohomology}
    H^2_{p}(X)(1) \hookrightarrow \End_{\Q_p}(W)
\end{equation}

\begin{lem}\label{Lemma:nilpotent operator on Kuga-Satake}
    Let $N_{\KS} \in \End_{K^{\ur}_0}(\Dpst(W))$
    be the monodromy operator on
$\Dpst(W)$.
Then
$
N_{\KS} = \rho^{\KS} (N_{p, 2})
$
where $N_{p,2} \in \so(H^2_{\pst}(X)(1))$ is the monodromy operator on $H^2_{\pst}(X)(1)$.
\end{lem}

\begin{proof}
Let $\{ s_{\alpha}\}_{\alpha}\subset W^{\otimes}$ be a set of tensors defining $\GSpin(H^2_p(X)(1)) \subset \GL_{\Q_p}(W)$.
Here $W^{\otimes}$ is the direct sum of all the $\Q_p$-vector spaces which can be constructed from $W$ by taking duals and tensor products.
Then
$\so(H^2_{p}(X)(1)) \oplus \Q_ph^{\deg} \subset \End_{\Q_p}(W)$
agrees with the subspace consisting of all endomorphisms that annihilate $\{ s_{\alpha}\}_\alpha$.

Since each $s_{\alpha}$ is $G_K$-invariant, it induces a tensor
$s_{\alpha, \pst} \in \Dpst(W)^\otimes$ which is annihilated by the monodromy operator.
Since
$\so(H^2_{\pst}(X)(1)) \oplus K^{\ur}_0h^{\deg} \subset \End_{K^{\ur}_0}(\Dpst(W))$
agrees with the subspace consisting of all endomorphisms that annihilate $\{ s_{\alpha, \pst}\}_\alpha$, we have
$
N_{\KS} \in \so(H^2_{\pst}(X)(1)) \oplus K^{\ur}_0h^{\deg}.
$
Since $N_{\KS}$ is nilpotent,
it follows that $N_{\KS} \in \so(H^2_{\pst}(X)(1))$.
The homomorphism
\[
H^2_{\pst}(X)(1) \to \End_{K^{\ur}_0}(\Dpst(W))
\]
induced by \eqref{Equation:Kuga-Satake embedding full cohomology} is compatible with monodromy operators, which implies that
$N_{\KS} =N_{p,2}$ in $\so(H^2_{\pst}(X)(1))$.
\end{proof}

\end{document}
